% ============================================================
% Formulario, argomento 16: vettori aleatori e vettori gaussiani (dispensa, capitolo 16)
% ============================================================
\argomento{Vettori aleatori e vettori gaussiani}{Dispensa: capitolo 16}[Vettori aleatori e gaussiani]

\voce{vettore}{Vettore aleatorio}{vettori sottolineati una volta, matrici due}
  {\vett{X} = [X_1, X_2, \ldots, X_n]^T \qquad f_{\vett{X}}(\vett{x}) = f_{X_1 \cdots X_n}(x_1, \ldots, x_n)}
  {\variante{CDF}{F_{\vett{X}}(\vett{x}) = P(X_1 \leq x_1, \ldots, X_n \leq x_n)}}

\voce{vettmedie}{Vettore delle medie}{la media si fa elemento per elemento}
  {\vett{\mu}_{\vett{X}} = E[\vett{X}] = [\mu_{X_1}, \mu_{X_2}, \ldots, \mu_{X_n}]^T}
  {}

\voce{matrcorr}{Matrice di correlazione}{simmetrica; diagonale: valori quadratici medi}
  {\matr{R}_{\vett{X}} = E\left[ \vett{X}\vett{X}^T \right], \qquad R_{X_i X_j} = E[X_i X_j]}
  {}

\voce{matrcov}{Matrice di covarianza}{simmetrica; diagonale: varianze}
  {\matr{C}_{\vett{X}} = E\left[ \left(\vett{X} - \vett{\mu}_{\vett{X}}\right) \left(\vett{X} - \vett{\mu}_{\vett{X}}\right)^T \right],
   \qquad C_{X_i X_j} = COV(X_i, X_j) = \rho_{X_i X_j}\, \sigma_{X_i} \sigma_{X_j}}
  {\variante{due variabili}{\matr{C}_{\vett{X}} = \begin{bmatrix} \sigma_1^2 & \rho\, \sigma_1 \sigma_2 \\
                                                                  \rho\, \sigma_1 \sigma_2 & \sigma_2^2 \end{bmatrix}}}

\voce{covcorr}{Covarianza dalla correlazione}{}
  {\matr{C}_{\vett{X}} = \matr{R}_{\vett{X}} - \vett{\mu}_{\vett{X}} \vett{\mu}_{\vett{X}}^T}
  {\variante{caratterizzazione sintetica}{\left(\vett{\mu}_{\vett{X}}, \matr{C}_{\vett{X}}\right): \ n \text{ medie e }
                                         \frac{n(n+1)}{2} \text{ elementi della matrice}}}

\voce{nonnegativa}{La matrice di covarianza è definita non negativa}{$\vett{q}$ vettore qualsiasi}
  {\vett{q}^T \matr{C}_{\vett{X}}\, \vett{q} = \mathrm{VAR}\left( \vett{q}^T \vett{X} \right) \geq 0}
  {\variante{due variabili}{\vett{q}^T \matr{C}_{\vett{X}} \vett{q} = q_1^2 \sigma_1^2 + q_2^2 \sigma_2^2 + 2 q_1 q_2\, \rho\, \sigma_1 \sigma_2}}

\voce{vettgauss}{Vettore gaussiano $\mathcal{N}(\vett{\mu}_{\vett{X}}, \matr{C}_{\vett{X}})$}{congiuntamente gaussiane, $\det \matr{C}_{\vett{X}} > 0$}
  {f_{\vett{X}}(\vett{x}) = \frac{1}{\sqrt{(2\pi)^n \det \matr{C}_{\vett{X}}}}\;
   \exp\left\{ -\frac{1}{2} \left( \vett{x} - \vett{\mu}_{\vett{X}} \right)^T \matr{C}_{\vett{X}}^{-1}
   \left( \vett{x} - \vett{\mu}_{\vett{X}} \right) \right\}}
  {\variante{componenti indipendenti}{\matr{C}_{\vett{X}} \text{ diagonale}, \quad f_{\vett{X}}(\vett{x}) = \prod_{i=1}^{n} f_{X_i}(x_i)}
   \variante{con $n = 1$ ridà}{\mathcal{N}(\mu_X, \sigma_X^2)}}

\voce{bivariata}{Gaussiana bivariata}{$\abs{\rho} < 1$}
  {f_{X_1 X_2}(x_1, x_2) = \frac{1}{2\pi \sigma_1 \sigma_2 \sqrt{1 - \rho^2}}\; \exp\left\{ -\frac{q(x_1, x_2)}{2(1 - \rho^2)} \right\}}
  {\variante{con}{q(x_1, x_2) = \frac{(x_1 - \mu_1)^2}{\sigma_1^2} - 2\rho\, \frac{(x_1 - \mu_1)(x_2 - \mu_2)}{\sigma_1 \sigma_2} +
                  \frac{(x_2 - \mu_2)^2}{\sigma_2^2}}
   \variante{determinante}{\det \matr{C}_{\vett{X}} = \sigma_1^2 \sigma_2^2 (1 - \rho^2)}
   \variante{inversa}{\matr{C}_{\vett{X}}^{-1} = \frac{1}{\sigma_1^2 \sigma_2^2 (1 - \rho^2)} \begin{bmatrix} \sigma_2^2 & -\rho\sigma_1\sigma_2 \\
     -\rho\sigma_1\sigma_2 & \sigma_1^2 \end{bmatrix}}
   \variante{curve di uguale densità}{\text{ellissi centrate in } (\mu_1, \mu_2)}}

\voce{cerchi}{Probabilità dei cerchi}{gaussiana bivariata standard con $\rho = 0$}
  {P\left( X_1^2 + X_2^2 \leq k^2 \right) = 1 - e^{-\frac{k^2}{2}}}
  {\variante{$k = 1, 2, 3$}{0{,}39, \quad 0{,}86, \quad 0{,}99 \quad (\text{e lo stesso per le ellissi, con qualsiasi } \rho)}}

\voce{gaussincorr}{Incorrelate e indipendenti}{per variabili congiuntamente gaussiane}
  {COV(X_i, X_j) = 0 \ \text{per } i \neq j \iff X_1, \ldots, X_n \text{ indipendenti}}
  {\variante{non basta che siano gaussiane una per una}{}}

\voce{gausscond}{Marginale e condizionata della bivariata}{caso standardizzato: medie 0, varianze 1}
  {X_1 \sim \mathcal{N}(0, 1) \qquad X_2 \mid X_1 = x_1 \;\sim\; \mathcal{N}\left(\rho\, x_1,\ 1 - \rho^2\right)}
  {\variante{media e varianza condizionate}{E[X_2 \mid X_1 = x_1] = \rho\, x_1, \quad \mathrm{VAR}(X_2 \mid X_1 = x_1) = 1 - \rho^2}}

\voce{proprietagauss}{Proprietà dei vettori gaussiani}{}
  {\text{bastano } \vett{\mu}_{\vett{X}} \text{ e } \matr{C}_{\vett{X}} \text{ a caratterizzarli completamente}}
  {\variante{trasformazioni lineari}{\text{restano gaussiani}}
   \variante{sottoinsiemi e componenti}{\text{gaussiani, con le loro medie e covarianze}}
   \variante{condizionamenti}{\text{gaussiani}}
   \variante{incorrelazione}{\text{equivale all'indipendenza}}}

\voce{trasflineare}{Trasformazione lineare di un vettore}{$\matr{A}$ è $m \times n$, $\vett{b}$ è $m \times 1$}
  {\vett{Z} = \matr{A}\vett{X} + \vett{b}: \qquad \vett{\mu}_{\vett{Z}} = \matr{A}\vett{\mu}_{\vett{X}} + \vett{b}
   \qquad \matr{C}_{\vett{Z}} = \matr{A}\, \matr{C}_{\vett{X}}\, \matr{A}^T}
  {\variante{le formule valgono per ogni vettore; se $\vett{X}$ è gaussiano}{\vett{Z} \sim
     \mathcal{N}(\vett{\mu}_{\vett{Z}}, \matr{C}_{\vett{Z}})}
   \variante{due variabili, $Z = a_1 X_1 + a_2 X_2$}{\sigma_Z^2 = a_1^2 \sigma_1^2 + a_2^2 \sigma_2^2 + 2 a_1 a_2\,
     \sigma_1 \sigma_2\, \rho_{12}}}

\voce{mediagauss}{Media di $n$ gaussiane indipendenti}{tutte $\mathcal{N}(\mu, \sigma^2)$}
  {\overline{X} = \frac{X_1 + \cdots + X_n}{n} \sim \mathcal{N}\left( \mu, \frac{\sigma^2}{n} \right)}
  {\variante{deviazione standard}{\frac{\sigma}{\sqrt{n}}}}

\voce{generacorr}{Costruire una coppia gaussiana correlata}{$W_1, W_2$ gaussiane standard indipendenti}
  {X_1 = W_1 \qquad X_2 = \rho\, W_1 + \sqrt{1 - \rho^2}\; W_2}
  {\variante{risultato}{\matr{C}_{\vett{X}} = \begin{bmatrix} 1 & \rho \\ \rho & 1 \end{bmatrix}}}
